\scanpage{4}
\uline{Lemme 2} Soit $L_*$ un $k$-module $\frac12$ simplicial, \deleted{\unsure{tel que}} \deleted{\unsure{(projectif)}} \deleted{flasque,} \deleted{que,} $q\geq q'\geq0$ deux entiers, $M$ un $k$-module, \inserted{\unsure{on suppose}, si $q=q'$, $\pi_{\unsure{0}}(L_*)=0$.} Alors
\[
\operatorname{Hom}_{\unsure{\widehat\Delta\mathrm{Ab}|\mathrm{Zab}}}
 (\underbrace{L_*\otimes_k\Lambda^q\Psi_{*,k}}_{\substack{\text{(on utilise seult}\\\text{que ce complexe est nul en dim }<q\text{)}}},\Lambda^{q'}\Psi_{*,k}\otimes_k M)=0.
\]


\deleted{Si $L'_*$ est le \unsure{complexe de chaînes associé}} \inserted{Le complexe de chaînes associé à} \deleted{\unsure{$L_*$, le complexe associé à}} $L_*\otimes_k\Lambda^q\Psi_{*,k}$ est \deleted{\unsure{un complexe}} \deleted{\unsure{de $k$-modules gradués, donc}} de la forme
\[
\ldots\longrightarrow L'_{q+1}\longrightarrow L'_q\longrightarrow0\longrightarrow0\ldots
\]
tandis que $\Lambda^{q'}\Psi_{*,k}\otimes_k M$ est $M$ placé en degré $q'$, \inserted{i.e. $M[q']$} si $q'<q$, il est évident que $\operatorname{Hom}_{\unsure{\bullet}}(L'_{\bullet},M_{\bullet})=0$, pour des raisons de \unsure{degrés|degré}. Si $q'=q$,
\diagram{\begin{tikzcd}[ampersand replacement=\&]
L'_{q+1}\arrow[r,"{\mathrm{surj}}"]\arrow[d] \& L'_q\arrow[r]\arrow[d] \&0\,\ldots\\
0\arrow[r]\& M\arrow[r]\&0
\end{tikzcd}}
alors $\pi_{\unsure{0|q}}(L_*)=0\ \Longrightarrow\ L'_{q+1}\xrightarrow{d}L'_q$ est surjectif, donc $L'_q\longrightarrow M$ est \unsure{nul|nulle} pour tout hom. de complexes \inserted{de chaînes} $L'_{q}\longrightarrow M[q]$, cqfd.

\uline{Théorème} Soit $n\geq0$. Il existe un homomorphisme de complexes de cochaînes dans $\unsure{\widehat{\Delta}_{\mathrm{ab}}}$
\diagram{\begin{tikzcd}[ampersand replacement=\&,column sep=tiny,font=\scriptsize]
0\arrow[r]\& \overset{0}{\Gamma^n\Phi_*}\arrow[r]\arrow[d,"u_0"]\&
\Gamma^{n-1}\Phi_*\otimes\Psi_*\arrow[r]\arrow[d,"u_1"]\&
\Gamma^{n-2}\Phi_*\otimes\Lambda^2\Psi_*\arrow[r]\arrow[d,"u_2"]\&
\cdots\arrow[r]\& \Lambda^n\Psi_*\arrow[r]\arrow[d,"u_n"]\&0\arrow[r]\arrow[d]\&0\\
0\arrow[r]\&\Phi_*\arrow[r,"{\wedge T}"']\&\Lambda^2\Phi_*\arrow[r,"{\wedge T}"']\&
\Lambda^3\Phi_*\arrow[r]\&\cdots\arrow[r,"{\wedge T}"']\&
\Lambda^{n+1}\Phi_*\arrow[r,"{\wedge T}"']\&\Lambda^{n+2}\Phi_*\arrow[r]\&\cdots
\end{tikzcd}}
(nb : on écrit $\Phi_*$ et $\Psi_*$ pour $\Phi_{*k}$, $\Psi_{*k}$, et $\otimes$ signifie $\otimes_k$), tel que $u_n$ soit l'isom. can. de $\Lambda^n\Psi_*$ sur $\operatorname{Ker}(\Lambda^{n+1}\Phi_*\twoheadrightarrow\Lambda^{n+1}\Psi_*\hookrightarrow\Lambda^{n+2}\Phi_*)=_{\mathrm{df}}Z(\Lambda^{\unsure{n+1}}\Phi_*)$
% archive sheet number: 3
